Traditional code synthesis relies on formal logic or rigid prompt engineering. We introduce \textbf{Vibe Coding}---a methodology where high-dimensional biological metaphors serve as executable functional specifications.

\subsection{Formal Definition of Vibe Coding}
\textbf{Definition 1 (Vibe Coding):}
Let $\mathcal{V}$ be a high-dimensional vibe space representing semantic intent, $\mathcal{M}$ a metaphor vocabulary derived from biological functional archetypes, and $\mathcal{C}$ an executable code space. Vibe Coding is defined as a mapping:
\begin{equation}
    \phi: \mathcal{V} \times \mathcal{M} \rightarrow \mathcal{C}
\end{equation}
realized by a large language model $f_{\theta}$ with parameters $\theta$:
\begin{equation}
    \phi(v, m) = f_{\theta}(\text{prompt}(v) \oplus \text{context}(m))
\end{equation}
subject to the functional constraint:
\begin{equation}
    \forall c \in \mathcal{C}, \; \text{Compile}(c) = \text{True} \land \text{TypeCheck}(c) = \text{True}
\end{equation}
The metaphor vocabulary $\mathcal{M}$ includes specialized mappings for agent roles:
\begin{itemize}
    \item \textbf{T-Cell (Audit)}: $\{scan, identify, neutralize\}$
    \item \textbf{B-Cell (Specialization)}: $\{verify, validate, synthesize\}$
    \item \textbf{Macrophage (Cleanup)}: $\{merge, seal, integrate\}$
\end{itemize}
A fundamental constraint is enforced: $\forall c \in \mathcal{C}$, $c$ must satisfy $\text{Compile}(c) = \text{True} \land \text{TypeCheck}(c) = \text{True}$.

\subsection{Complexity Analysis of Algorithm 1}
The time complexity of a single generation in the Dialectical Protocol is dominated by the adversarial audit phase. Let $K$ be the number of agents in the committee and $L$ the average context length. The audit phase scales as $O(K \cdot L^2)$ assuming standard transformer attention mechanisms. However, since the Digital Twin sandbox operates in parallel, the total wall-clock time is $O(L^2 + T_{val})$, where $T_{val}$ is the constant-time test suite execution.

\subsection{Vibe Coding vs. Traditional Prompting}
Unlike traditional prompt engineering which specifies \textit{how} a computation should proceed, Vibe Coding specifies the \textit{biological objective}, aligning with the recent paradigm of intent-based programming \cite{karpathy2025vibecoding, sorensen2024vibe}. For example, a Vibe input of ``Ligate the Bayesian specialist into the loop'' implies a structural integration that maintains context-identity while requiring Merkle-root (SHA-512) validation for state-integrity. This allows for the autonomous discovery of domain-specific scaling laws \cite{lewis2022scaling} without manual parameter tuning.

\subsection{Implementation Stack}
The Vibe Coding substrate utilizes a heterogeneous model distribution to ensure audit diversity:
\begin{itemize}
    \item \textbf{Auditors}: Mistral-Large-2 (Semantic Audit) and Gemini 1.5 Pro (Context Audit).
    \item \textbf{Executors}: Phi-3 and Neural-Chat for low-latency mutation generation.
    \item \textbf{Sandbox}: Docker-based isolation with Intel oneAPI acceleration.
\end{itemize}
